% Computer Science A --- Java LTS releases (17 vs.\ 21) and sealed classes.
% \texttt{cs-web-adv/} module; companion to \texttt{web-adv.tex}.

\runningheader{Computer Science A}{Java LTS \& sealed types}{Page \thepage\ of \numpages}

\begingradingrange{csajavalts}

\uplevel{\par\textbf{Java platform (17 vs.\ 21)}\par}

\question[6]
Compare Java 17 and Java 21. Why would a new project choose Java 21?
\begin{solutionorbox}[5.5in]
\textbf{Definition:} Java 17 (LTS, September 2021) and Java 21 (LTS, September 2023) are both Long-Term Support releases. Java 21 finalizes major concurrency and data-manipulation features that were preview-only or absent in Java 17.

\textbf{Purpose:} Teams pick Java 21 for higher-throughput server workloads, cleaner pattern-matching code, and a more uniform Collections API---without abandoning the stability expectations of an LTS line.

\textbf{Key differences:}
\begin{itemize}
  \item \textbf{Concurrency:} Java 17 maps each \lstinline[style=javainline]|Thread| to a heavy OS thread (large stack, kernel context switches). A typical JVM supports only a few thousand active platform threads before memory pressure dominates. Java 21 introduces \textbf{Virtual Threads} (Project Loom): lightweight threads scheduled by the JVM onto a small pool of carrier threads. When a virtual thread blocks on I/O, the JVM parks it and reuses the carrier for another task---yielding async-style throughput with synchronous-looking code.
  \item \textbf{Pattern matching:} Java 17 standardizes \lstinline[style=javainline]|instanceof| pattern matching (type check and bind in one step). Java 21 finalizes \textbf{pattern matching for \lstinline[style=javainline]|switch|} and \textbf{record patterns}, so cases can match object types, deconstruct records, apply \lstinline[style=javainline]|when| guards, and handle \lstinline[style=javainline]|null| directly.
  \item \textbf{Sealed types + exhaustiveness:} Sealed classes arrived in Java 17; Java 21 pairs them with exhaustive \lstinline[style=javainline]|switch| expressions so the compiler rejects missing cases when every permitted subclass is handled.
  \item \textbf{Sequenced Collections (JEP 431):} Before Java 21, first/last element access differed by collection type (\lstinline[style=javainline]|.get(0)|, \lstinline[style=javainline]|.iterator().next()|, \lstinline[style=javainline]|.get(size()-1)|). Java 21 adds \lstinline[style=javainline]|SequencedCollection|, \lstinline[style=javainline]|SequencedSet|, and \lstinline[style=javainline]|SequencedMap| with uniform \lstinline[style=javainline]|getFirst()|, \lstinline[style=javainline]|getLast()|, \lstinline[style=javainline]|reversed()|, and related methods.
  \item \textbf{Garbage collection:} Java 21 improves generational ZGC out of the box, reducing allocation stalls under load compared with manual ZGC tuning common on Java 17 deployments.
\end{itemize}

\textbf{Example (Java 17 --- \lstinline[style=javainline]|instanceof| pattern matching):}
\begin{lstlisting}[style=java]
Object obj = "Hello from Java 17";
if (obj instanceof String s) {
    System.out.println(s.toUpperCase()); // no (String) cast
}
\end{lstlisting}

\textbf{Example (Java 21 --- virtual threads):}
\begin{lstlisting}[style=java]
try (var executor = Executors.newVirtualThreadPerTaskExecutor()) {
    IntStream.range(0, 100_000).forEach(i ->
        executor.submit(() -> {
            Thread.sleep(Duration.ofSeconds(1));
            return i;
        })
    );
}
\end{lstlisting}

\textbf{Example (Java 21 --- record patterns in \lstinline[style=javainline]|switch|):}
\begin{lstlisting}[style=java]
record Point(int x, int y) {}

Object obj = new Point(10, 20);
String result = switch (obj) {
    case Integer i -> "Integer: " + i;
    case String s  -> "String: " + s;
    case Point(int x, int y) -> "Point at " + x + ", " + y;
    default -> "Unknown";
};
\end{lstlisting}

\textbf{Recommendation:} For new server-side or full-stack Java projects, Java 21 is the stronger default because virtual threads simplify high-concurrency web services and pattern matching reduces boilerplate around domain data.
\end{solutionorbox}

\newpage

\question[6]
What are sealed classes, and how do they support a functional style in modern Java?
\begin{solutionorbox}[4.5in]
\textbf{Definition:} A \textbf{sealed class} or \textbf{sealed interface} restricts which types may extend or implement it. The base type declares a finite \lstinline[style=javainline]|permits| list; each permitted subclass must be \lstinline[style=javainline]|final|, \lstinline[style=javainline]|sealed|, or \lstinline[style=javainline]|non-sealed|.

\textbf{Purpose:} Sealed types model \textbf{algebraic data types} (ADTs) from functional programming: a value can only be one of a known set of variants. Combined with Java 21 pattern matching, the compiler enforces \textbf{exhaustiveness}---every permitted case must be handled, so runtime \lstinline[style=javainline]|IllegalArgumentException| guards become unnecessary.

\textbf{Clarification:} Sealed classes were \emph{finalized in Java 17}. Java 21 finalizes pattern matching for \lstinline[style=javainline]|switch|, which is the feature that unlocks compile-time exhaustiveness checking on sealed hierarchies.

\textbf{Example (sealed hierarchy):}
\begin{lstlisting}[style=java]
public sealed interface Shape permits Circle, Rectangle, Square {}

public final class Circle implements Shape {
    private final double radius;
    public Circle(double radius) { this.radius = radius; }
    public double radius() { return radius; }
}
// Rectangle and Square similarly declared final
\end{lstlisting}

\textbf{Example (exhaustive area calculation without a \lstinline[style=javainline]|default| case):}
\begin{lstlisting}[style=java]
public double getArea(Shape shape) {
    return switch (shape) {
        case Circle c    -> Math.PI * c.radius() * c.radius();
        case Rectangle r -> r.length() * r.width();
        case Square s    -> s.side() * s.side();
        // no default: compiler verifies all permitted shapes are covered
    };
}
\end{lstlisting}

\textbf{Benefits:}
\begin{itemize}
  \item \textbf{Domain control:} Third-party code cannot silently add new subclasses that bypass your API rules.
  \item \textbf{Compiler safety:} Adding a new permitted subclass breaks the build anywhere a non-exhaustive \lstinline[style=javainline]|switch| exists, surfacing missing logic before runtime.
  \item \textbf{Cleaner code:} Eliminates defensive \lstinline[style=javainline]|else| branches and ``unknown type'' exceptions at the bottom of handlers.
\end{itemize}
\end{solutionorbox}

\endgradingrange{csajavalts}
